% ---------------------------------------------------------------------------
% Appendix: the five SKILL.md files, verbatim.
%
% Upload this file AND the skills/ folder next to it, then \input this once.
% The \lstinputlisting paths resolve from the PROJECT ROOT, not from wherever this
% file sits, so skills/ has to be at the top level of the Overleaf project even if
% this .tex lives in a subfolder.
% The .md files are pulled in with \lstinputlisting rather than pasted, so the
% thesis can never hold a copy that has drifted from the repo. Refresh them with
%   python3 scripts/collect_skill_files.py
%
% Needs only listings + xcolor, both already loaded.
%
% Two things this style is working around, both real:
%   - the preamble renames \lstlistingname to "Algorithmus", so every listing
%     here uses title= rather than caption=; a captioned listing would print
%     "Algorithmus N" in an English thesis and appear in the list of listings,
%     where a 173-line skill file does not belong.
%   - the files contain 50 em dashes, 9 right arrows and 2 en dashes, which
%     listings cannot typeset under the deprecated utf8x inputenc. The literate
%     mapping below handles them inside the style, leaving the preamble alone.
% breaklines and postbreak are not optional: the longest line in
% secure-coding-v1 is 525 columns.
% ---------------------------------------------------------------------------

\lstdefinestyle{skillmd}{
  basicstyle=\ttfamily\footnotesize,
  breaklines=true, columns=fullflexible, keepspaces=true,
  showstringspaces=false, frame=single, framesep=4pt,
  rulecolor=\color{black!30}, xleftmargin=12pt,
  aboveskip=1.2em, belowskip=0.6em,
  postbreak=\mbox{\textcolor{gray}{$\hookrightarrow$}\space},
  literate={—}{{-{}-{}-}}3 {–}{{-{}-}}2 {→}{{$\rightarrow$}}1
}

\section{The skill files}
\label{app:skills}

The five \texttt{SKILL.md} files used in this thesis are reproduced verbatim
below. The three merge-conflict versions are the successive outputs of the design
loop described in Section~\ref{sec:loop_procedure}; the two benchmark skills were
each written in a single version. Where a line is too long for the page it is
broken and marked with $\hookrightarrow$; that arrow is typesetting, not part of
the file.

\subsection{merge-conflict-resolve-v1}

The control, 277 words: a four-way choice between keeping either side, combining
them, and writing a new resolution, with no worked examples and no edge cases.

\lstinputlisting[style=skillmd]{skills/merge-conflict-resolve-v1.md}

\subsection{merge-conflict-resolve-v2}

742 words. Adds the ordered pattern hierarchy and the output discipline v1
lacked, including the cap of $|a| + |b|$ characters.

\lstinputlisting[style=skillmd]{skills/merge-conflict-resolve-v2.md}

\subsection{merge-conflict-resolve-v2.1}

1102 words. Moves the output discipline ahead of the pattern hierarchy, replaces
the character cap with content rules and a post-hoc self-check, and hoists the
one-side-empty rule into the pattern test itself.

\lstinputlisting[style=skillmd]{skills/merge-conflict-resolve-v2.1.md}

\subsection{secure-coding-v1}

The SALLM skill, 1677 words. Written in a single version: it inherits the
conclusions of the merge-conflict loop without having been iterated on itself.

\lstinputlisting[style=skillmd]{skills/secure-coding-v1.md}

\subsection{swebench-repair-v1}

The SWE-bench Lite skill, 501 words, single-version on the same basis.

\lstinputlisting[style=skillmd]{skills/swebench-repair-v1.md}
